\section{Syzygy, dark planes and certificates}\label{sec:dark}

Define
\[
P_A=\{d_1=d_2,\ d_3=d_4\},\qquad
P_B=\{d_1=d_4,\ d_3=d_2\}.
\]

\begin{theorem}[Pointwise dark-plane decomposition]\label{thm:darkplanes}
Over every integral domain,
\begin{equation}\label{eq:darkpoints}
L(d)=Q(d)=0\quad\Longleftrightarrow\quad d\in P_A\cup P_B.
\end{equation}
On $P_A$ the original signed involution persists, and on $P_B$ the
alternative signed involution persists.  Consequently every selected moment,
and hence the selected amplitude, vanishes on the union.
\end{theorem}
\begin{proof}
Eliminate $d_1$ using $L=0$, so $d_1=d_2-d_3+d_4$.  Then
\[
Q=d_2d_4-(d_2-d_3+d_4)d_3=(d_2-d_3)(d_4-d_3).
\]
The absence of zero divisors gives either
$d_2=d_3$, $d_1=d_4$, which is $P_B$, or
$d_4=d_3$, $d_1=d_2$, which is $P_A$.  The converse follows by
substitution.  On each plane the matching diagonal perturbation commutes with
the stated signed involution, so the argument of
\cref{prop:baseline-dark} applies at every perturbation value.
\end{proof}

\noindent\textbf{Formal status.}
The equivalence \cref{eq:darkpoints} is \evidence{KERNEL-CHECKED} as
\path{M3A.ell_quad_zero_iff_planes}; Lean proves it over every integral domain.
The persistent commutators and the all-depth consequences are checked by
\path{M3A.rOrig_commutes_hamiltonian_of_planeA},
\path{M3A.rAlt_commutes_hamiltonian_of_planeB} and
\path{M3A.moments_dark_of_linear_quadratic_zero}.  M8 then gives actual
amplitude darkness.

\begin{proposition}[Reduced ideal decomposition]\label{prop:darkideal}
Over any field of characteristic different from $2$,
\begin{equation}\label{eq:darkdecomp}
(L,Q)=(d_1-d_2,d_3-d_4)\cap(d_1-d_4,d_3-d_2).
\end{equation}
Consequently the exact-dark scheme is the reduced union of two
two-dimensional linear planes.
\end{proposition}
\begin{proof}
The factorisation in the proof of \cref{thm:darkplanes} gives the two distinct
prime components after eliminating $L$.  Pulling them back to the original
polynomial ring gives \cref{eq:darkdecomp}; their intersection is radical.
\end{proof}

This scheme-theoretic ideal equality is \evidence{PROVED} in the paper but is
outside the current Lean scope.  Separating it from
\cref{thm:darkplanes} prevents the kernel-checked pointwise theorem from being
silently upgraded to an unformalised statement about polynomial ideals.

On a fixed ray $v=(a,b,c,d)$, the two planes read
\begin{equation}
P_A:\ a=b,\ c=d,\qquad P_B:\ a=d,\ c=b.
\end{equation}
They intersect in the scalar line $a=b=c=d$, an equality also checked in Lean
by \path{M3A.onPlaneA_and_onPlaneB_iff_scalarLine}.

\begin{figure}[ht]
\centering\begin{tikzpicture}[>=Latex,scale=.9,
  plane/.style={draw,thick,fill opacity=.24,text opacity=1},
  lab/.style={font=\small,align=center}]
\draw[->] (-4,0)--(4.4,0) node[right] {$\ell$ direction};
\draw[->] (0,-2.5)--(0,2.8) node[above] {$Q$ direction inside $\ell=0$};
\fill[blue!12] (-.75,-2.2) rectangle (.75,2.2);
\draw[blue!60!black,thick] (-.75,-2.2) rectangle (.75,2.2);
\draw[plane,fill=green!50] (-.6,-1.8)--(.55,-.9)--(.6,1.8)--(-.55,.9)--cycle;
\draw[plane,fill=orange!60] (-.55,-.9)--(.6,-1.8)--(.55,.9)--(-.6,1.8)--cycle;
\node[lab] at (2.4,1.4) {$\ell\ne0$\\linear opening};
\node[lab] at (0,2.45) {$\ell=0,\ Q\ne0$:\\quadratic opening};
\node[lab] at (0,-.05) {two dark planes\\$P_A\cup P_B$};
\end{tikzpicture}

\caption{Schematic direction stratification.  The ambient hyperplane
\(\ell=0\) contains the two exact-dark planes; points elsewhere on the
hyperplane have a quadratic opening.}
\label{fig:darkstrata}
\end{figure}

\subsection{Certificate discipline}

A signed involution $S$ with $SH_0=H_0S$, $Se_0=e_0$ and
$Se_5=-e_5$ is sufficient for baseline darkness.  Under perturbation the
commutator \([S,D(d)]\) may become non-zero.  This proves only that this
certificate has failed.  Exact opening requires a surviving coefficient of
$N$, of the moment series, or of $K$.  Conversely, a point of $V(L,Q)$
may be dark while failing to preserve a chosen presentation of $S$.

\begin{proposition}[Coefficient syzygy]\label{prop:syzygy}
The triple \((L,Q,R)\) satisfies the exact relation
\[
d_2d_4L+(d_2+d_4)Q+R=0.
\]
The relation explains why a nominal cubic numerator coefficient supplies no
third stratum after $L=Q=0$.
\end{proposition}

\begin{proof}
Expand the left-hand side and cancel the four cubic monomials.  This is
\cref{eq:Rsyzygy} rewritten and is also \evidence{KERNEL-CHECKED} as
\path{M3A.coefficient_syzygy}.
\end{proof}

This syzygy is the algebraic reason that the fixed-ray classification has only
linear, quadratic and exact-zero branches.  It is not a statement that the
cubic moment itself is generated by the two earlier moments; the tagged depth
information is addressed in \cref{sec:gf1}.
